% =====================================================================
% paper_insertions.tex — LaTeX blocks to paste into the SOCR AI Bot
% Data Descriptor (Overleaf). Each block is marked with WHERE it goes.
% British spelling used throughout to match the existing manuscript
% (tumour / normalisation / optimiser / artefact / modelling).
%
% INTEGRITY NOTES (read first):
%  - All numbers are from Alex's real runs. No FID is claimed (not computed).
%  - Do NOT add Alex's row to Table 3 (tab:metric_scores): its PSNR/SSIM use a
%    per-slice paired protocol and are not comparable to the class-paired numbers
%    there. A separate results paragraph is provided instead.
%  - Coordinate with the lead author (A. Shankar) before inserting.
%  - Values marked [VERIFY] / [TBD] must be filled from the current logs/code.
% =====================================================================


% ---------------------------------------------------------------------
% INSERT 1 — DATA RECORDS: replace the \textcolor{red}{...} block
% (finalised prose; adds Alex's model subfolder, flagged for release status)
% ---------------------------------------------------------------------
The repository contains the following top-level folders.

\paragraph{\texttt{/models/}} Pre-trained model weight files in PyTorch \texttt{.pt} or
\texttt{.pth} format, one subfolder per model:
\begin{itemize}
  \item \texttt{braingen\_GAN\_TCGA\_v1/}
  \item \texttt{braingen\_cGAN\_BraTS\_v1/}
  \item \texttt{braingen\_WaveletGAN\_v1/}
  \item \texttt{braingen\_Diffuser\_v1/}
  \item \texttt{braingen\_CondDiffuser\_BraTS\_v1/}  % [VERIFY release status]
  \item \texttt{braingen\_gan3d\_BraTS\_32\_v1/}
  \item \texttt{braingen\_gan3d\_BraTS\_64\_v1/}
\end{itemize}
Each subfolder contains the generator and (where applicable) discriminator weights, a JSON
metadata file describing the architecture and the conditioning vocabulary, and a Python
script demonstrating sample generation.

\paragraph{\texttt{/samples/}} Pre-generated synthetic image samples organised by model.
For the 2D models, 1{,}000 samples per conditional class are provided as 16-bit PNG files at
the model's native resolution. For each multi-contrast 2D model, the per-channel outputs
(T1, T1ce, T2, FLAIR, and the tumour segmentation mask) are saved as separate files sharing a
common filename stem. For the 3D models, 200 sample volumes per model are provided as
compressed NIfTI files (\texttt{.nii.gz}) with isotropic spacing. File names encode the
conditioning class and sample index (for example, \texttt{ax\_mid\_tumour\_0042\_t1.png}).

\paragraph{\texttt{/metadata/}} A CSV file (\texttt{samples\_metadata.csv}) describing every
synthetic image in \texttt{/samples/}. Columns include: \texttt{model}, \texttt{sample\_id},
\texttt{anatomical\_plane} (axial/sagittal/coronal), \texttt{location}
(superior/middle/inferior, left/middle/right, or anterior/middle/posterior),
\texttt{tumour\_present} (0/1), \texttt{contrast} (T1/T1ce/T2/FLAIR/mask),
\texttt{resolution\_px}, \texttt{file\_path}, and \texttt{file\_format}.

\paragraph{\texttt{/metrics/}} Numerical evaluation results in CSV format: per-model FID,
SSIM, MSE, and PSNR with bootstrap 95\% confidence intervals on the held-out test partition.
A separate file (\texttt{splits.csv}) lists the patient-level train/validation/test
assignments used for both datasets.

\paragraph{\texttt{/code/}} A mirror of the public source-code repository (also available on
GitHub, see Code Availability) containing the model definitions, training scripts, evaluation
pipeline, and the Shiny web application source.

A complete tree of file paths, sizes, and MD5 checksums is provided as
\texttt{file\_manifest.csv} at the repository root.


% ---------------------------------------------------------------------
% INSERT 2 — USAGE NOTES: replace the \textcolor{red}{To Finalize ...} block
% ---------------------------------------------------------------------
For programmatic reuse, the released model weights can be downloaded from the data repository
(see Data Records and Data Availability) and loaded through the provided PyTorch inference
scripts. Each model subfolder includes a self-contained example that instantiates the network,
loads the weights, accepts the conditioning arguments documented in the accompanying
\texttt{README}, and writes generated images (and, where applicable, the paired segmentation
mask) to disk. The README files specify the conditioning vocabulary, the expected input
intensity ranges, and the output tensor layout for each model. The Shiny application source is
included in the \texttt{/code/} folder for users who wish to deploy a private instance of the
interface behind their own compute.


% ---------------------------------------------------------------------
% INSERT 3 — TABLE 1 (tab:sota_comparison): new row in the
% "Models released with the SOCR AI Bot" block (after braingen_Diffuser_v1)
% ---------------------------------------------------------------------
braingen CondDiffuser BraTS\_v1 & 2025 & Conditional DDPM/DDIM ($\epsilon$-pred, U-Net) & $256\!\times\!256$ (2D) & BraTS 2023 GLI ($n\!=\!1{,}252$) & FLAIR (cond.\ on T1+mask+atlas) & PSNR, SSIM, bgSNR, loss-floor ablation & Yes \\


% ---------------------------------------------------------------------
% INSERT 4 — TABLE 2 (tab:training_config): new row (ALL VALUES FILLED)
% Deployed model = job 52208104 (full data): 34 epochs, 30 h wall-clock
% (State=TIMEOUT at the --time=30h limit; per-epoch checkpoints kept).
% Params = 85.3 M (build_model() -> 85.261185 M).
% NOTE: trained on an NVIDIA RTX PRO 6000 (96 GB), unlike the V100/Titan XP rows.
% ---------------------------------------------------------------------
braingen\_CondDiffuser\_BraTS\_v1 & 1e--4 & AdamW & 32 & 34 & $\sim$30 & 85.3 & 89 \\


% ---------------------------------------------------------------------
% INSERT 5 — METHODS: new \paragraph after braingen_Diffuser_v1, in
% "Two-dimensional generative models"
% ---------------------------------------------------------------------
\paragraph{braingen\_CondDiffuser\_BraTS\_v1.} A conditional denoising diffusion probabilistic
model\cite{ho2020denoising} with $\epsilon$-prediction, trained to synthesise FLAIR slices
conditioned on a subject's own anatomy. Whereas braingen\_Diffuser\_v1 is unconditional, this
model receives an eight-channel spatial conditioning tensor---the subject's co-registered
T1-weighted slice, the tumour segmentation mask, and six probabilistic lobe maps from the
ICBM452 atlas\cite{mazziotta2001probabilistic} warped into the SRI24\cite{rohlfing2010sri24}
reference space---concatenated channel-wise with the noised FLAIR input to a U-Net denoiser
(channel widths $128\!\to\!256\!\to\!384\!\to\!512$, self-attention at the coarsest scale). The
atlas was registered once into SRI24 space with ANTs symmetric normalisation
(SyN)\cite{avants2008syn} and reused across subjects, exploiting the fact that BraTS volumes
are already mutually co-registered; per-subject anatomy therefore enters only through the T1
and mask channels. The forward and reverse diffusion processes follow
\cref{eq:diff_forward,eq:diff_reverse,eq:diff_mean}: the FLAIR target is scaled to $[-1,1]$
while the conditioning channels are retained in $[0,1]$, and the network minimises the
noise-prediction objective
$\mathbb{E}_{t,x_{0},\epsilon}\lVert \epsilon - \epsilon_{\theta}(x_{t},t,c)\rVert_{2}^{2}$,
where $c$ is the conditioning tensor. Optimisation used AdamW (learning rate $10^{-4}$), batch
size 32, mixed-precision, and an exponential moving average (decay 0.999) of the weights for
sampling. Slices were intensity-normalised per subject by the $[1,99]$ percentile of brain
voxels and zero-padded from $240\!\times\!240$ to $256\!\times\!256$ (a power of two for the
U-Net) rather than resampled. Generation uses 200 deterministic DDIM\cite{song2021ddim} steps.


% ---------------------------------------------------------------------
% INSERT 6 — TECHNICAL VALIDATION: new \paragraph in "Qualitative inspection"
% (or as its own short subsection). Leads with the informational-floor finding.
% ---------------------------------------------------------------------
\paragraph{Anatomical conditioning and the noise-prediction floor.} On the BraTS test
partition, braingen\_CondDiffuser\_BraTS\_v1 reproduces subject-specific anatomy (ventricular
morphology, cortical folding, grey/white-matter contrast) and tumour location while retaining
stochastic texture diversity across independent samples of the same conditioning
(\cref{fig:conddiff_results}). An ablation on the conditioning tensor isolates the value of the
anatomical prior: with only the tumour mask and atlas as conditioning, the $\epsilon$-prediction
training loss plateaus at $\approx 6\times10^{-3}$ irrespective of model capacity or training
length, whereas adding the subject's own T1 channel lowers the plateau to
$\approx 4.2\times10^{-3}$ (a $\sim$30\% reduction). This is consistent with an
information-theoretic reading in which the attainable noise-prediction error is lower-bounded by
the conditional variance of the target given the conditioning: weakly-determining conditioning
imposes an irreducible floor that stronger, subject-specific conditioning relaxes. Pixelwise
fidelity to the paired real slice, computed over the brain region on 100 held-out slices (one
stochastic sample each, 200 DDIM steps), was PSNR $17.1\pm2.8$\,dB and SSIM $0.40\pm0.13$, and
the generated backgrounds were essentially noise-free (background signal-to-noise ratio
$66.6\pm35.5$). Because each output is a single stochastic draw that deliberately does not copy
the reference, these full-reference scores are expected to fall below those of deterministic
translation models and are reported as conservative fidelity bounds rather than as directly
comparable to \cref{tab:metric_scores}.


% ---------------------------------------------------------------------
% INSERT 7 — FIGURE (Technical Validation). Copy Alex's grid into Overleaf's
% Figures/ as conddiff_eval_grid.png  (source: eval_grid_20260626_133819.png)
% ---------------------------------------------------------------------
\begin{figure}[h]
    \centering
    \includegraphics[width=0.9\columnwidth]{Figures/conddiff_eval_grid.png}
    \caption{Samples from \textit{braingen\_CondDiffuser\_BraTS\_v1}. Leftmost panel: a
    held-out real FLAIR slice; remaining panels: independent synthetic samples generated from
    the same anatomical conditioning (subject T1 + tumour mask + atlas). Per-panel PSNR, SSIM,
    and background SNR are annotated. All samples share the subject's gross anatomy and tumour
    location while differing in fine texture, illustrating conditioned generation rather than
    reconstruction.}
    \label{fig:conddiff_results}
\end{figure}


% ---------------------------------------------------------------------
% INSERT 8 — AUTHOR CONTRIBUTIONS: replace the A.L. sentence
% ---------------------------------------------------------------------
R.K.\ contributed to the conditional GAN implementations and the web application development.
A.L.\ designed and implemented the anatomy-conditioned diffusion model
(braingen\_CondDiffuser\_BraTS\_v1), including the atlas-to-SRI24 registration, the
conditioning pipeline, and the associated training and evaluation experiments.


% ---------------------------------------------------------------------
% INSERT 9 — references.bib: add these three entries (new cites used above)
% ---------------------------------------------------------------------
@article{rohlfing2010sri24,
  author  = {Rohlfing, Torsten and Zahr, Natalie M. and Sullivan, Edith V. and Pfefferbaum, Adolf},
  title   = {The {SRI24} multichannel atlas of normal adult human brain structure},
  journal = {Human Brain Mapping},
  volume  = {31}, number = {5}, pages = {798--819}, year = {2010},
  doi     = {10.1002/hbm.20906}
}

@article{mazziotta2001probabilistic,
  author  = {Mazziotta, John and Toga, Arthur and Evans, Alan and others},
  title   = {A probabilistic atlas and reference system for the human brain: International Consortium for Brain Mapping ({ICBM})},
  journal = {Philosophical Transactions of the Royal Society B},
  volume  = {356}, number = {1412}, pages = {1293--1322}, year = {2001},
  doi     = {10.1098/rstb.2001.0915}
}

@article{avants2008syn,
  author  = {Avants, Brian B. and Epstein, Charles L. and Grossman, Murray and Gee, James C.},
  title   = {Symmetric diffeomorphic image registration with cross-correlation: Evaluating automated labeling of elderly and neurodegenerative brain},
  journal = {Medical Image Analysis},
  volume  = {12}, number = {1}, pages = {26--41}, year = {2008},
  doi     = {10.1016/j.media.2007.06.004}
}


% ---------------------------------------------------------------------
% INSERT 10 (OPTIONAL) — ABSTRACT: extend the model enumeration.
% Replace "...a denoising diffusion implicit model, and two 3D conditional GANs"
% with:
% ---------------------------------------------------------------------
% ...an unconditional denoising diffusion implicit model, an anatomy-conditioned
% diffusion model that synthesises FLAIR conditioned on a subject's T1, tumour mask,
% and a probabilistic lobe atlas, and two 3D conditional GANs...


% ---------------------------------------------------------------------
% EXISTING BUGS SPOTTED (fix in the current manuscript, not additions):
%  (a) Methods cites \cite{song2021maximum} twice, but the bib key is song2021ddim.
%      -> change song2021maximum to song2021ddim (or add the alias to references.bib).
%  (b) Table 2 (tab:training_config): column spec is {lcccccccc} = 9 columns but
%      every row has 8 fields -> a trailing empty column. Drop one 'c' for cleanliness.
% ---------------------------------------------------------------------
